% Source: https://artofproblemsolving.com/community/c3h3192620_2021_bmmt_individual_round__berkley_mini_math_tournament
% Stable identity: aops-topic-3192620-q19
Unfold any sequence of crossed faces into the plane. A shortest path does not revisit a face, since the straight chord inside that convex face shortens such a revisit. It also cannot pass through a vertex: the total face angle there is 240 degrees, so one side of the path has angle less than 180 degrees and can be shortened. Thus it suffices to check simple paths in the dual cube, of which there are 18. The following exact certificate unfolds all of them and retains only straight segments crossing each shared edge in order and in its interior. Coordinates $(u,v)$ in the code represent $(u,v\sqrt3)$, so squared lengths are $u^2+3v^2$.\begin{verbatim}
from fractions import Fraction as F
from itertools import product
faces=list(product([0,1],[2,3],[4,5]))
start,goal=faces[0],faces[-1]
def sub(a,b):return (a[0]-b[0],a[1]-b[1])
def dot(a,b):return a[0]*b[0]+3*a[1]*b[1]
def cross(a,b):return a[0]*b[1]-a[1]*b[0]
def reflect(p,a,b):
 v=sub(b,a);t=dot(v,sub(p,a))/dot(v,v)
 return (2*(a[0]+t*v[0])-p[0],2*(a[1]+t*v[1])-p[1])
P=(F(1),F(1,3)); lengths=[]; path_count=0
def search(face,coords,used,edges):
 global path_count
 if face==goal:
  path_count+=1
  Q=tuple(sum(coords[v][j] for v in face)/3 for j in [0,1])
  vec=sub(Q,P);last=0
  for a,b in edges:
   e=sub(b,a);w=sub(a,P);den=cross(vec,e)
   if not den:return
   t=cross(w,e)/den;u=cross(w,vec)/den
   if not(last<t<1 and 0<u<1):return
   last=t
  lengths.append(dot(vec,vec));return
 for nxt in faces:
  shared=set(face)&set(nxt)
  if len(shared)!=2 or nxt in used:continue
  a,b=[coords[v] for v in shared]
  old=next(v for v in face if v not in shared)
  new=next(v for v in nxt if v not in shared)
  cc={v:coords[v] for v in shared}
  cc[new]=reflect(coords[old],a,b)
  search(nxt,cc,used|{nxt},edges+[(a,b)])
search(start,{0:(F(0),F(0)),2:(F(2),F(0)),4:(F(1),F(1))},{start},[])
assert path_count==18 and min(lengths)==F(28,3)
\end{verbatim}
The minimum squared distance is $\dfrac{28}{3}$, so $\boxed{m+n=31}$. One attaining unfolded chain has endpoint centres $(1,\dfrac{\sqrt3}{3})$ and $(3,5\dfrac{\sqrt3}{3})$, whose squared separation is $4+\dfrac{16}{3}=\dfrac{28}{3}$.
\begin{asy}
usepackage("amsmath");
size(7cm);
pair A=(0,0),B=(2,0),C=(1,sqrt(3)),Ap=(3,sqrt(3)),Bp=(2,2*sqrt(3)),Cp=(4,2*sqrt(3));draw(A--B--C--cycle);draw(B--Ap--C);draw(C--Bp--Ap);draw(Ap--Cp--Bp);pair P=(1,sqrt(3)/3),Q=(3,5*sqrt(3)/3);draw(P--Q,blue+linewidth(1.5));dot(P);dot(Q);label("$P$",P,W);label("$Q$",Q,E);label("$2$",(A+B)/2,S);
\end{asy}
